\documentclass{situationist-preprint}

\title{Defining \enquote{Ought} in Terms of Reasons}
\msid{MTE-11-b}
\msfamily{MTE-11}
\msdiscipline{Metaethics}
\mstype{Consequence}
\mskeywords{ought; reasons; reasons first; objective and evidence-relative ought; miners case; agent-relativity}

\begin{document}
\maketitle

\begin{abstract}
MTE-11-a defines \enquote{ought} as what the agent has strongest reason to choose among her options. We derive what this definition implies for the varieties of ought. Relativizing the reasons to a base yields the objective and the evidence-relative ought as instances of one definition, and the definition resolves the miners case without positing two senses of the word. It further follows that ought is agent-relative wherever reasons are, that ought is gappy where permission is not, and that the fact that one ought to act is not a further reason to act.
\end{abstract}

\section{The definition and its relativization}

Definition~3.3 of MTE-11-a says that, in a situation with choice set $S$, the agent ought to choose $a$ if $a\in S$ and the reasons for $a$ are stronger than the reasons for every other option in $S$. Definition~3.7 relativizes permissibility to a \emph{base}, which assigns to each situation the body of considerations from which the reasons are drawn. Since ought is unique permission (Proposition~4.4 of MTE-11-a), the same relativization applies to ought.

\begin{definition}\label{d:based}
Relative to a base $B$, the agent \emph{ought$_B$} to choose $a$ in $S$ if $a\in S$ and the reasons that the considerations in $B$ would provide are stronger for $a$ than for every other option in $S$.
\end{definition}

The \emph{objective} ought takes as base all the facts of the situation; the \emph{evidence-relative} ought takes the facts the agent possesses; Parfit's \emph{belief-relative} ought takes the facts her beliefs represent \citep{Parfit2011}. Each is an instance of Definition~\ref{d:based}, and each inherits the structure proved in Section~4 of MTE-11-a, since those proofs use only the ordering of total strengths and not the base from which the strengths are drawn.

\section{Consequences}

\paragraph{The miners.} Ten miners are trapped in one of two shafts, the agent does not know which, and flood water is coming. Blocking the right shaft saves all ten, blocking the wrong one saves none, and blocking neither saves nine \citep{KolodnyMacFarlane2010}. Relative to the full base, the agent ought to block the shaft the miners are in. Relative to her evidence, which gives equal support to each shaft, the reasons for blocking either shaft are those of an even chance of saving ten, and the reasons for blocking neither are those of saving nine for certain; she ought$_{\mathrm{ev}}$ to block neither. Both verdicts are true, and they do not conflict, since they are oughts relative to different bases. The case is often taken to show that \enquote{ought} is ambiguous or context-sensitive in a way that resists a single analysis. On the reasons-first definition it has a single analysis with a parameter, and the deliberative ought, which the agent asks about when she asks what to do, is the ought relative to the reasons she possesses.

\paragraph{Agent-relativity.} The reason relation of Assumption~3.1 of MTE-11-a holds between a fact, an agent, a response and a situation. A fact may be a reason for one agent and not for another: that a child is mine is a reason for me to save her that is not, in the same way, a reason for a stranger. By Definition~3.3, ought inherits this relativity. The same holds of value-for-an-agent, which Definition~3.6 of MTE-11-a, applied to the reasons an agent has to prefer one outcome to another, defines from possibly agent-relative reasons. This is why the agent-relative value functions of ETH-02, which represent constraints and prerogatives, are not an artifice: they are the value-for-an-agent that the reasons themselves define.

\paragraph{Gaps.} If two options tie in strength and none is stronger, Definition~3.3 yields no ought, while Proposition~4.3 of MTE-11-a guarantees that some option is permissible. Ought is gappy where permission is not. This is the correct result: in a choice between two equally good ways of helping, it is false that the agent ought to take the first and false that she ought to take the second, and true that she may take either.

\paragraph{No double counting.} The fact that one ought to choose $a$ is not a further reason to choose it. If it were, the strength of the reasons for $a$ would depend on whether $a$ is supported by the strongest reasons, and Definition~3.3 would be circular. The definition therefore entails that the deliberative force of an ought is exhausted by the reasons it summarizes. An agent who acts because she ought to acts for the reasons that make it the case that she ought, or acts on a belief that there are such reasons.

\section{Objection and reply}

\begin{objection}
\citet{Broome2013} takes ought as the basic normative notion, a relation between an agent and a proposition, and defines reasons from it. The many senses of \enquote{ought} that the miners case brings out are then evidence that ought is not a single weighing of reasons, since the weighing cannot be the same in the objective and the deliberative case.
\end{objection}

\begin{reply}
The weighing is the same; what differs is the base from which the reasons are drawn. Definition~\ref{d:based} gives one relation of strongest reason with a parameter, and the miners case fixes different values of the parameter. The reverse order of definition, from ought to reasons, is the one that Theorem~5.6 of MTE-11-a shows cannot succeed, since two situations can agree on every deontic fact and differ in their reasons; Objection~1 of MTE-11-a takes up Broome's version of it, which adds the notion of explanation to the deontic base.
\end{reply}

\section{Conclusion}

Defining ought as strongest reason, and relativizing reasons to a base, yields the objective, evidence-relative and belief-relative oughts as one notion with a parameter. The definition resolves the miners case, makes ought agent-relative wherever reasons are, leaves ought gappy where permission is not, and denies that an ought adds to the reasons it summarizes. These are consequences of the primacy of reasons established in MTE-11-a, not further assumptions.

\begin{thebibliography}{9}\small
\bibitem[Broome(2013)]{Broome2013} Broome, J. (2013). \emph{Rationality Through Reasoning}. Wiley-Blackwell.
\bibitem[Kolodny and MacFarlane(2010)]{KolodnyMacFarlane2010} Kolodny, N. and MacFarlane, J. (2010). Ifs and oughts. \emph{The Journal of Philosophy}, 107(3), 115--143.
\bibitem[Parfit(2011)]{Parfit2011} Parfit, D. (2011). \emph{On What Matters}, vol.~1. Oxford University Press.
\end{thebibliography}

\end{document}
